% =====================================================================
\chapter{2 次元とドープ系}
\label{ch:2d}
% =====================================================================

1 次元では、低い結合次元の MPS が局所観測量に十分な prior になる。2 次元の幅 $W$ のシリンダーでは、MPS で状態を表すのに必要な結合次元が幅について指数的に増える($\chi\sim e^{cW}$)。この章では、MPS prior が高くつく 2 次元で、MPS を使わない平均場 prior がどこまで代わりになるかを調べる。

\section{設定}
\begin{itemize}
  \item 系:$W=4$ と $W=6$ のシリンダー($L_x=8$、円周方向は周期、軸方向は開放)。64 量子ビットと 96 量子ビット。$U=4,8$。
  \item 参照状態:DMRG の基底状態(結合次元 \todo{W=4: 192、W=6: 512 を確認})。2 次元の DMRG の状態は完全には収束しておらず、副格子の磁化を持つ(スピン回転の対称性を破る)。
  \item prior:UHF、UHF-sym、参照状態を $\chi$ に切り詰めた MPS。
  \item 観測量:中央のサイトの周りのオンサイト $ZZ$、隣接と距離 2 の $Z_\up Z_\up$、二重占有、$S^zS^z$、$y$ 方向のホッピング。
  \item 利得は分散の厳密な式による($n_m=100$)。
\end{itemize}

\section{$W=4$ シリンダー}
\label{sec:w4}

\begin{table}[t]
\centering\small
\caption{$W=4$ シリンダー、$U=4$ の利得(分散の厳密な式、$n_m=100$)。}
\label{tab:w4}
\begin{tabular}{@{}lcccc@{}}
\toprule
観測量 & UHF & UHF-sym & 切断 $\chi{=}16$ & 切断 $\chi{=}64$ \\
\midrule
オンサイト $ZZ$ & 27.0 & 27.0 & 13.0 & 28.6 \\
二重占有 & 6.9 & 3.4 & 9.7 & 23.8 \\
ホッピング($y$) & 12.9 & 12.9 & 16.5 & 16.5 \\
$Z_\up Z_\up$($x$ 方向の隣) & 0.97 & 8.68 & 6.8 & 8.9 \\
$S^zS^z$($x$ 方向の隣) & 0.88 & 6.59 & 4.5 & 7.8 \\
\bottomrule
\end{tabular}
\end{table}

表\ref{tab:w4}から次のことが分かる。
\begin{itemize}
  \item \textbf{オンサイト $ZZ$ とホッピングでは、UHF が $\chi=64$ の MPS に並ぶ。} $U=8$ ではオンサイト $ZZ$ で UHF 141 対 $\chi=64$ 143(天井 145)である。2 次元では MPS の忠実度が急に下がる($U=4$、$\chi=8$ で $F=0.055$。1 次元 $L=16$ では同じ $\chi$ で 0.95)ので、計算量 $O(N^3)$ の平均場がこれに並ぶことには実用上の意味がある。
  \item \textbf{スピン相関では UHF は損をし、UHF-sym で $\chi=64$ に近い水準まで直る。}
  \item \textbf{二重占有では UHF-sym の方が悪い}(6.9 → 3.4、$U=8$ では 10.5 → 5.1)。第\ref{sec:sumobs}節のとおり利得は項ごとの外れで決まる。参照状態の DMRG がスピン回転の対称性を破って $\ev{Z_{i\up}}\ne0$ を持つので、項ごとには共線の UHF の方が合う。これは参照状態の性質によるもので、実験の状態が対称性を保つなら逆になるはずである。
\end{itemize}

\paragraph{モンテカルロの値との違い.}以前の表は反復 50 回のモンテカルロの値で、オンサイト $ZZ$ が 44.0(UHF)・44.5($\chi=64$)と、天井 28.9 を大きく超えていた。同じ測定データで全 prior を評価したため、標準の推定量の分散の偶然の大きさが全列に共通にかかった。同じ乱数の種で再実行すると同じ外れ方を再現したので、計算の誤りではない。

\section{$W=6$ シリンダー}
\label{sec:w6}
\todo{crm\_2d\_w6\_results\_U4p0.tsv / U8p0.tsv の G\_exact 列で表を作る。以前の要約の「UHF $G=157$ 対 $\chi=32$ の $G=96$」(モンテカルロ、反復 30 回)を置き換える}

\section{ドープ系}
\label{sec:doped}
\todo{crm\_2d\_doped\_results\_U4p0/U8p0.tsv(W=4、ホール 4 個)と crm\_2d\_w6\_doped\_results\_*.tsv(W=6、ホール 6 個、窓の位置 x0=0, 2)の G\_exact で表を作る}

$1/8$ のホールを入れると、平均場は電荷の川(ストライプ)を自発的に作るが、その電荷の変調が真の状態より鋭すぎる。サイトごとの密度 $\ev{n_i}$ では UHF の外れが $0.13$〜$0.16$ に達し、平均場 prior はこの観測量では役に立たない。\todo{数値は再計算の結果で確認する}

\section{まとめ}
\todo{再計算の結果がそろってから書く}
